%%%%%%%%%%%%%%%%%%%%%%%%%%%%%%%%%%%%%%%%%
% Lachaise Assignment
% LaTeX Template
% Version 1.0 (26/6/2018)
%
% This template originates from:
% http://www.LaTeXTemplates.com
%
% Authors:
% Marion Lachaise & François Févotte
% Vel (vel@LaTeXTemplates.com)
%
% License:
% CC BY-NC-SA 3.0 (http://creativecommons.org/licenses/by-nc-sa/3.0/)
% 
%%%%%%%%%%%%%%%%%%%%%%%%%%%%%%%%%%%%%%%%%

%----------------------------------------------------------------------------------------
%	PACKAGES AND OTHER DOCUMENT CONFIGURATIONS
%----------------------------------------------------------------------------------------

\documentclass{article}
\usepackage{amsthm}
\usepackage{listings}
\usepackage{xcolor}
\usepackage{graphicx}
\usepackage{cite}
\usepackage{longtable}
\usepackage{array}
\usepackage{booktabs}

\lstset{
	language=Python,
	basicstyle=\ttfamily\small,
	keywordstyle=\color{blue},
	stringstyle=\color{red},
	commentstyle=\color{purple},
	showstringspaces=false,
	numbers=left,
	numberstyle=\tiny\color{gray},
	breaklines=true,
	frame=single,
	captionpos=b
}
\newtheorem{definition}{Definition}[section]

%%%%%%%%%%%%%%%%%%%%%%%%%%%%%%%%%%%%%%%%%
% Lachaise Assignment
% Structure Specification File
% Version 1.0 (26/6/2018)
%
% This template originates from:
% http://www.LaTeXTemplates.com
%
% Authors:
% Marion Lachaise & François Févotte
% Vel (vel@LaTeXTemplates.com)
%
% License:
% CC BY-NC-SA 3.0 (http://creativecommons.org/licenses/by-nc-sa/3.0/)
% 
%%%%%%%%%%%%%%%%%%%%%%%%%%%%%%%%%%%%%%%%%

%----------------------------------------------------------------------------------------
%	PACKAGES AND OTHER DOCUMENT CONFIGURATIONS
%----------------------------------------------------------------------------------------

\usepackage{amsmath,amsfonts,stmaryrd,amssymb} % Math packages

\usepackage{enumerate} % Custom item numbers for enumerations

\usepackage[ruled]{algorithm2e} % Algorithms

\usepackage[framemethod=tikz]{mdframed} % Allows defining custom boxed/framed environments

\usepackage{listings} % File listings, with syntax highlighting
\lstset{
	basicstyle=\ttfamily, % Typeset listings in monospace font
}

%----------------------------------------------------------------------------------------
%	DOCUMENT MARGINS
%----------------------------------------------------------------------------------------

\usepackage{geometry} % Required for adjusting page dimensions and margins

\geometry{
	paper=a4paper, % Paper size, change to letterpaper for US letter size
	top=2.5cm, % Top margin
	bottom=3cm, % Bottom margin
	left=2.5cm, % Left margin
	right=2.5cm, % Right margin
	headheight=14pt, % Header height
	footskip=1.5cm, % Space from the bottom margin to the baseline of the footer
	headsep=1.2cm, % Space from the top margin to the baseline of the header
	%showframe, % Uncomment to show how the type block is set on the page
}

%----------------------------------------------------------------------------------------
%	FONTS
%----------------------------------------------------------------------------------------

\usepackage[utf8]{inputenc} % Required for inputting international characters
\usepackage[T1]{fontenc} % Output font encoding for international characters

\usepackage{XCharter} % Use the XCharter fonts

%----------------------------------------------------------------------------------------
%	COMMAND LINE ENVIRONMENT
%----------------------------------------------------------------------------------------

% Usage:
% \begin{commandline}
%	\begin{verbatim}
%		$ ls
%		
%		Applications	Desktop	...
%	\end{verbatim}
% \end{commandline}

\mdfdefinestyle{commandline}{
	leftmargin=10pt,
	rightmargin=10pt,
	innerleftmargin=15pt,
	middlelinecolor=black!50!white,
	middlelinewidth=2pt,
	frametitlerule=false,
	backgroundcolor=black!5!white,
	frametitle={Command Line},
	frametitlefont={\normalfont\sffamily\color{white}\hspace{-1em}},
	frametitlebackgroundcolor=black!50!white,
	nobreak,
}

% Define a custom environment for command-line snapshots
\newenvironment{commandline}{
	\medskip
	\begin{mdframed}[style=commandline]
}{
	\end{mdframed}
	\medskip
}

%----------------------------------------------------------------------------------------
%	FILE CONTENTS ENVIRONMENT
%----------------------------------------------------------------------------------------

% Usage:
% \begin{file}[optional filename, defaults to "File"]
%	File contents, for example, with a listings environment
% \end{file}

\mdfdefinestyle{file}{
	innertopmargin=1.6\baselineskip,
	innerbottommargin=0.8\baselineskip,
	topline=false, bottomline=false,
	leftline=false, rightline=false,
	leftmargin=2cm,
	rightmargin=2cm,
	singleextra={%
		\draw[fill=black!10!white](P)++(0,-1.2em)rectangle(P-|O);
		\node[anchor=north west]
		at(P-|O){\ttfamily\mdfilename};
		%
		\def\l{3em}
		\draw(O-|P)++(-\l,0)--++(\l,\l)--(P)--(P-|O)--(O)--cycle;
		\draw(O-|P)++(-\l,0)--++(0,\l)--++(\l,0);
	},
	nobreak,
}

% Define a custom environment for file contents
\newenvironment{file}[1][File]{ % Set the default filename to "File"
	\medskip
	\newcommand{\mdfilename}{#1}
	\begin{mdframed}[style=file]
}{
	\end{mdframed}
	\medskip
}

%----------------------------------------------------------------------------------------
%	NUMBERED QUESTIONS ENVIRONMENT
%----------------------------------------------------------------------------------------

% Usage:
% \begin{question}[optional title]
%	Question contents
% \end{question}

\mdfdefinestyle{question}{
	innertopmargin=1.2\baselineskip,
	innerbottommargin=0.8\baselineskip,
	roundcorner=5pt,
	nobreak,
	singleextra={%
		\draw(P-|O)node[xshift=1em,anchor=west,fill=white,draw,rounded corners=5pt]{%
		Question \theQuestion\questionTitle};
	},
}

\newcounter{Question} % Stores the current question number that gets iterated with each new question

% Define a custom environment for numbered questions
\newenvironment{question}[1][\unskip]{
	\bigskip
	\stepcounter{Question}
	\newcommand{\questionTitle}{~#1}
	\begin{mdframed}[style=question]
}{
	\end{mdframed}
	\medskip
}

%----------------------------------------------------------------------------------------
%	WARNING TEXT ENVIRONMENT
%----------------------------------------------------------------------------------------

% Usage:
% \begin{warn}[optional title, defaults to "Warning:"]
%	Contents
% \end{warn}

\mdfdefinestyle{warning}{
	topline=false, bottomline=false,
	leftline=false, rightline=false,
	nobreak,
	singleextra={%
		\draw(P-|O)++(-0.5em,0)node(tmp1){};
		\draw(P-|O)++(0.5em,0)node(tmp2){};
		\fill[black,rotate around={45:(P-|O)}](tmp1)rectangle(tmp2);
		\node at(P-|O){\color{white}\scriptsize\bf !};
		\draw[very thick](P-|O)++(0,-1em)--(O);%--(O-|P);
	}
}

% Define a custom environment for warning text
\newenvironment{warn}[1][Warning:]{ % Set the default warning to "Warning:"
	\medskip
	\begin{mdframed}[style=warning]
		\noindent{\textbf{#1}}
}{
	\end{mdframed}
}

%----------------------------------------------------------------------------------------
%	INFORMATION ENVIRONMENT
%----------------------------------------------------------------------------------------

% Usage:
% \begin{info}[optional title, defaults to "Info:"]
% 	contents
% 	\end{info}

\mdfdefinestyle{info}{%
	topline=false, bottomline=false,
	leftline=false, rightline=false,
	nobreak,
	singleextra={%
		\fill[black](P-|O)circle[radius=0.4em];
		\node at(P-|O){\color{white}\scriptsize\bf i};
		\draw[very thick](P-|O)++(0,-0.8em)--(O);%--(O-|P);
	}
}

% Define a custom environment for information
\newenvironment{info}[1][Info:]{ % Set the default title to "Info:"
	\medskip
	\begin{mdframed}[style=info]
		\noindent{\textbf{#1}}
}{
	\end{mdframed}
}
 % Include the file specifying the document structure and custom commands

% --- Ambiente para citar o parecer (definição local; não faz parte do template) ---
\newenvironment{parecer}
{\begin{quote}\itshape\color{gray!70!black}\small}
	{\end{quote}}

% --- Rótulo de identificação do revisor (definição local; não faz parte do template) ---
\newcommand{\revisor}[1]{\textbf{[Revisor #1]}}

%----------------------------------------------------------------------------------------
%	ASSIGNMENT INFORMATION
%----------------------------------------------------------------------------------------


\title{IA Neurossimbólica para Perguntas e Respostas sobre Modelos Metabólicos em Escala Genômica: um Estudo Comparativo entre RAG Vetorial, GraphRAG e OG-RAG \\ \large Carta de Resposta aos Pareceres (Rebuttal)} % Title of the assignment

\author{Giovani Tavares\\ \texttt{giovanitavares@usp.br}} % Author name and email address

\date{Universidade de São Paulo --- \today} % University, school and/or department name(s) and a date

%----------------------------------------------------------------------------------------

\begin{document}
	
	\maketitle % Print the title
	
	%----------------------------------------------------------------------------------------
	%	INTRODUÇÃO
	%----------------------------------------------------------------------------------------
	
	\section{Introdução}
	
	Agradeço aos revisores pela avaliação criteriosa e pelos comentários
	construtivos. Respondo abaixo, ponto a ponto, apenas aos critérios que
	receberam classificação \emph{Razoavelmente} ou inferior, indicando as
	revisões incorporadas à versão final da proposta. As demais observações
	positivas foram registradas com apreço e não requerem réplica.
	
	% =====================================================================
	\section{Critério 3 --- Metodologia: baselines e métricas}
	% =====================================================================
	
	\subsection{Resumo das objeções}
	\begin{parecer}
		\revisor{2} ``Diferencial não tem métricas bem definidas'';
		``Ausência de baselines e métricas''.
	\end{parecer}
	
	\subsection{Resposta}
	
	O conjunto de métricas RAGAS (\textit{faithfulness}, \textit{answer relevance},
	\textit{context precision} e \textit{context recall}) citadas na Seção~2.2 serão
	utilizadas como instrumento de avaliação. Elas serão calculadas sobre as três
	arquiteturas apresentadas, entre as quais o RAG Vetorial será o \textit{baseline}.
	Dessa forma, a proposta será alterada de modo a deixar as métricas descritas em
	uma subseção específica, ao invés de no meio da bibliografia, e o RAG Vetorial
	será explicitado como o \textit{baseline}.
	
	Cabe ainda esclarecer que o desenho experimental já contempla, por construção,
	o isolamento do efeito da camada simbólica: as três arquiteturas (RAG Vetorial,
	GraphRAG e OG-RAG) compartilham o mesmo LLM gerador, o mesmo corpus-fonte e o
	mesmo conjunto de perguntas de avaliação, variando exclusivamente o mecanismo
	de recuperação. Assim, o RAG Vetorial (sem qualquer estrutura simbólica, isto
	é, sem ontologia) funciona como \textit{baseline} contra o qual GraphRAG e OG-RAG
	são comparados, e as diferenças observadas nas métricas do RAGAS podem ser
	atribuídas à presença da camada simbólica.
	
	\subsection{Alterações na proposta}
	
	As definições formais das métricas foram movidas da Seção~2.2 (Bibliografia)
	para uma subseção exclusiva dentro da metodologia --- ``Definição das Métricas
	de Avaliação'' ---, permanecendo na Bibliografia apenas a citação ao RAGAS com
	remissão à nova subseção. Adicionalmente, o papel do RAG Vetorial como
	\textit{baseline} sem camada simbólica foi explicitado no objetivo específico
	correspondente e na descrição da Fase~A do desenho experimental.
	
	% =====================================================================
	\section{Critério 6 --- Plano de Gestão de Dados}
	% =====================================================================
	
	\subsection{Resumo das objeções}
	\begin{parecer}
		\revisor{1} ``Não deixou muito explícito os tipos de dados e nem as
		formas de armazenamento''.\\
		\revisor{2} ``É necessário reformular com plano de gestão de dados''.\\
		\revisor{3} ``Descreve os tipos de dados, mas deve se atentar às formas de
		armazenamento, preservação e compartilhamento''.
	\end{parecer}
	
	\subsection{Resposta}
	
	Este foi o ponto convergente dos três pareceres e recebeu atenção prioritária
	na revisão. A proposta passou a contar com uma seção dedicada de Plano de Gestão
	de Dados, organizada em três eixos:
	
	\begin{description}
		\item[Tipos de dados.] O projeto produz e manipula: (i) dados-fonte de domínio
		do modelo \textit{E. coli core} do repositório BiGG Models, em formatos SBML e
		JSON; (ii) a ontologia OntoGEM adaptada, em formato OWL; (iii) o corpus textual
		comum em \textit{markdown}, gerado deterministicamente a partir dos dados
		estruturados; (iv) o conjunto anotado de perguntas de competência, respostas de
		referência e classificações por tipo de explicação, em formato tabular
		(CSV/JSON); (v) os artefatos intermediários das execuções (contextos recuperados
		e respostas geradas por cada arquitetura); e (vi) as métricas RAGAS calculadas e
		os resultados das análises estatísticas.
		
		\item[Armazenamento e preservação.] Durante a execução, dados e código serão
		versionados em repositório públicos do Github. Para preservação de longo
		prazo, os artefatos publicamente compartilháveis serão depositados em repositório
		público com atribuição de DOI (por exemplo, Zenodo), priorizando-se formatos
		abertos e não proprietários (SBML, OWL, JSON, CSV, \textit{markdown}).
		
		\item[Compartilhamento e licenciamento.] O conjunto de perguntas e respostas
		anotado, o código dos três \textit{pipelines} e os \textit{scripts} de análise
		serão disponibilizados publicamente sob licença de código aberto (por exemplo,
		MIT para código e CC-BY para dados), seguindo os princípios FAIR. O reuso de
		dados de terceiros (BiGG Models e OntoGEM) respeitará as licenças originais, com
		a devida atribuição. Não há manipulação de dados sensíveis ou pessoais, pois todo
		o material é de natureza técnico-científica e derivado de fontes públicas.
	\end{description}
	
	\subsection{Alterações na proposta}
	
	Foi inserida a seção ``Plano de Gestão de Dados'', contendo as subseções
	``Tipos de Dados Produzidos'', ``Armazenamento e Preservação'' e
	``Compartilhamento e Licenciamento'', conforme descrito acima.
	
	% =====================================================================
	\section{Critério 8 --- Liderança científica (ou potencial)}
	% =====================================================================
	
	\subsection{Resumo das objeções}
	\begin{parecer}
		\revisor{1} ``Não há evidências \ldots sobre ele ter atuado como líder
		em pesquisa ou algum potencial de liderança''.
	\end{parecer}
	
	\subsection{Resposta}
	
	Trata-se de uma proposta de mestrado e, portanto, de um pesquisador em início de
	carreira, estágio em que o critério deve ser lido em termos de \emph{potencial}
	de liderança, e não de um histórico consolidado. Nesse sentido, cabe registrar
	que a atuação profissional do proponente como \textit{Machine Learning Engineer}, tendo também concluído uma iniciação científica na área de grandes modelos de linguagem ao final de sua graduação como cientista da computação na mesma instituiação.
	
	% =====================================================================
	\section{Critério 9 --- Experiência na formação de pesquisadores}
	% =====================================================================
	
	\subsection{Resumo das objeções}
	\begin{parecer}
		\revisor{1} ``Não há descrição \ldots sobre experiência na formação de
		pesquisadores compatível com o estágio da carreira''.\\
		\revisor{3} ``Não há evidências \ldots mas o autor demonstra competência e
		domínio da pesquisa''.
	\end{parecer}
	
	\subsection{Resposta}
	
	Como o critério prevê expressamente, a experiência na formação de pesquisadores
	deve ser ``compatível com o estágio da carreira''. Em um projeto de mestrado, não
	se espera que o proponente já tenha orientado formalmente outros pesquisadores.
	
	% =====================================================================
	\section{Critério 10 --- Experiência / colaboração internacional}
	% =====================================================================
	
	\subsection{Resumo das objeções}
	\begin{parecer}
		\revisor{1} ``Não há descrição \ldots sobre experiências internacionais em
		pesquisas''.\\
		\revisor{3} ``Não há conhecimento do histórico acadêmico \ldots internacional
		\ldots mas o autor demonstra competência e domínio para tal''.
	\end{parecer}
	
	\subsection{Resposta}
	
	Este critério, em sua formulação, refere-se a ``experiência internacional em
	pesquisa após o doutoramento ou participação ativa em redes internacionais de
	colaboração'', condição naturalmente não aplicável a um proponente em nível de
	mestrado. O compromisso	com a publicação em código aberto e a disponibilização do conjunto de dados
	anotado, mencionados no Plano de Gestão de Dados, criam condições concretas para
	a inserção do proponente em redes internacionais de colaboração. A versão revisada
	da proposta torna esse enquadramento mais explícito.
	
	% =====================================================================
	\section{Considerações finais}
	% =====================================================================
	
	Reitero meu agradecimento pela avaliação. As revisões acima --- em especial a
	inclusão de um Plano de Gestão de Dados completo e o detalhamento das métricas e
	do \textit{baseline} --- foram integralmente incorporadas à versão revisada da
	proposta, que acredito atender de forma mais robusta aos critérios da FAPESP.
	
	\vspace{1.5em}
	\noindent Atenciosamente,\\[2em]
	\noindent\rule{6cm}{0.4pt}\\
	Giovani Tavares\\
	IME-USP
	
\end{document}